\section{Introduction}

In most African countries, the livestock sector supports the livelihoods of millions of rural households and plays an important role in poverty reduction. In livestock-rich countries such as Kenya, the sector contributes substantially to the national economy while providing income, employment, and food security across both rural and urban populations. Beyond primary production, livestock also supports employment throughout value chains, including dairy, meat, and leather processing industries. The livestock sector plays a major role in Kenya's food system, contributing approximately 12\% of the country's Gross Domestic Product (GDP), 40\% of agricultural GDP, and 22\% of the food system GDP \cite{bahtaLivestockSectorTransformation2023}.

Kenya's population is projected to reach approximately 96 million by 2050, with nearly half of the population residing in urban areas. Consequently, demand for animal-source foods is expected to increase substantially owing to population growth, urbanization, rising incomes, and changing dietary preferences \cite{bahtaLivestockSectorTransformation2023}. Meeting this growing demand will require sustained improvements in livestock productivity, efficient management of feed and water resources, and evidence-based planning to support the sustainable development of the livestock sector. Reliable forecasting frameworks that project livestock populations, production outputs, feed demand, and resource requirements are therefore essential for evidence-based planning, strategic investment, and sustainable livestock sector development.

Statistical time-series models are commonly used to forecast livestock populations using historical data. These methods estimate future livestock numbers based on past population trends and have been applied to support livestock planning and management. \cite{bahtaLivestockSectorTransformation2023} used autoregressive and probabilistic time-series models to forecast pig populations, while Ozbek compared autoregressive integrated moving average (ARIMA) and artificial neural network (ANN) models for forecasting cattle populations in Türkiye. Similarly, \cite{bahtaLivestockSectorTransformation2023} applied exponential smoothing to forecast beef cattle populations in Indonesia. Although these methods have demonstrated good predictive performance, they rely primarily on historical trends and do not explicitly represent biological processes such as reproduction, mortality, maturation, culling, and offtake. Consequently, their ability to evaluate the effects of biological parameters and management interventions on future livestock population dynamics is limited.

In contrast to statistical forecasting approaches, biological simulation models explicitly represent the demographic processes governing livestock populations. Animals are grouped into cohorts according to age, sex, or production stage, with transitions driven by biological processes such as reproduction, maturation, mortality, culling, and replacement. These models offer a biologically realistic representation of herd dynamics and enable the evaluation of alternative management strategies. \cite{bahtaLivestockSectorTransformation2023} developed an age- and sex-structured cattle population model that incorporates management practices and animal movements to simulate herd dynamics. Such models serve as valuable decision-support tools by linking biological processes to herd structure and productivity, they are generally developed for specific livestock species or production systems and do not explicitly integrate livestock population dynamics, production, environmental carrying capacity, and feed demand within a unified forecasting framework. Consequently, there remains a lack of reproducible analytical frameworks that simultaneously integrate livestock population dynamics, production, environmental carrying capacity, and feed demand to support long-term livestock sector planning.

To address these limitations, this study develops a cohort-based quarterly forecasting framework for Kenya's nine principal livestock value chains. The framework explicitly models biological cohort transitions through reproduction, mortality, maturation, culling, and offtake while linking livestock population dynamics to species-specific production, environmental carrying capacity, and feed demand across agro-ecological zones. By integrating these components within a single modelling framework, the proposed approach provides a comprehensive decision-support tool for forecasting livestock populations, production, and resource requirements, thereby supporting evidence-based livestock planning and policy development in Kenya.

